\section{Results: Gaps between ASB 018's goals (envisioned audit practices) and the audit practices it enables}
\label{sec:results}
From our analysis of study compliance with ASB 018, we found that all the audits we examined could be interpreted as compliant with our permissive interpretation of ASB 018's requirements language. 

In this section, we present our findings on the gaps between ASB 018's envisioned audit practices and the audit practices it enables, and how the design of ASB 018 enables these gaps. First, we present our findings on how ASB 018 envisions the \textit{purpose} of audits, since we use this to inform our interpretation of ASB 018's goals for each of the five audit stages. We organize the remainder our findings by the five audit stages outlined in Section~\ref{sec:methods}. In each subsection, we first present each of the goal-practice gaps and then discuss aspects of ASB 018's design that enable these gaps. Figure~\ref{fig:results-summary} summarizes our findings on the goal-practice gaps for each audit stage.

\textbf{Audit Purpose:} Audits help laboratories \textbf{(i) establish boundaries on software use that delineate the range of DNA samples for which the audit has demonstrated acceptable PGS system performance.} These boundaries, in turn, \textbf{(ii) shape protocols for software use in casework} and \textbf{(iii) provide assurance to stakeholders that the likelihood ratio evidence produced using the PGS system in an individual criminal case is reliable.} ASB 018 explicitly states this goal when defining internal validation as the ``acquisition of test data within the laboratory to [...] determin[e] [...] limitations of the system'' ~\cite[p.2]{asb018}. The factsheet provides additional context suggesting that these limitations are boundaries on software use in casework: validation studies ``define limitations for its use'' that ``establish the range of DNA profiles upon which the program may be used effectively''~\cite{asb018factsheet}. Through these statements, we understand  limitations as delineations between the types of DNA profiles for which internal validation studies have demonstrated acceptable PGS system performance that others may ``have confidence in''~\cite{asb018factsheet}, and, the types of DNA profiles for which the software has not been shown to be sufficiently reliable and should therefore not be used. 

\subsection{Audit Scope: Determining the Audit Target and Selecting Issues to Prioritize Investigating in the Audit.}
\label{sec:scope}
\subsubsection{Gap: Broader sociotechnical system vs. Probabilistic genotyping software.}
\label{sec:scope-gap1-target}

\textbf{\\Goal: Compliant audits evaluate the broader sociotechnical system that takes a crime scene sample as input and produces a likelihood ratio (LR) as output. This system encompasses both the PGS and the PGS user.} When using the PGS to produce a LR, a forensic lab analyst specifies key inputs to the software such as their estimate for the \textit{number of contributors (NoC)} to the DNA sample. ASB 018 highlights the ``user input parameters'' as key components of the broader sociotechnical system that audits must ``evaluat[e]''~\cite[R4.1.4]{asb018} and ``address''~\cite[Annex R4.1.4]{asb018}. Testing this broader sociotechnical system is also crucial to ASB 018's end goal for audits to provide assurance about the reliability of LR evidence---a \textit{system} output---presented in a criminal case.


\textbf{Practices:} \textbf{Several audits did not use the user's estimate of the \textit{number of contributors (NoC)} parameter: a key user input that significantly impacts the reliability of the final LR output and is frequently estimated incorrectly.} While the OCME and CBI labs' audit reports describe using user estimates of the NoC when testing software sensitivity and specificity, the PBSO and MSP labs' audit reports describe using the \textit{ground truth} NoC. These audits' use of the ground truth NoC is especially concerning given the frequency of incorrect NoC estimates in the audits that used user NoC estimates and their consequent impacts on the LR outputs. For instance, the OCME v2.7 audit report reveals that users underestimated the NOC for 25 out of 40 five-person samples and noted that several false negative results ``occurred with [five-person DNA mixtures] where the apparent NoC was [underestimated]''~\cite[p.21]{ocmeIV2.7}. 

\textbf{Some audits claimed that software users would be able to prevent observed software failures in casework, despite having no experiments that tested software users' ability to identify and prevent such failures.} For example, the OCME v2.7 audit observed several false negative LRs but claimed that a ``trained analyst'' would be able to notice that the software's intermediate outputs were ``unintuitive,'' re-run the software with a correction, and produce correct results (i.e., true positive LRs)~\cite[p.8]{ocmeIV2.7}. However, without incorporating the analyst's inspection of these intermediate outputs into the testing process, the audit provides no assurance that a user would correctly identify and fix such failures in casework. In other words, the audit report claims that the \textit{system} would produce a correct output when analyzing similar samples in casework but does not conduct tests that would provide evidence to support the claim and provide assurance of system reliability in casework.

\textbf{We also found evidence in several reports that suggested that the PGS developers, and not the labs, conducted parts or all of the audit.} This suggests that any PGS user decisions that were used in the audit may have come from the developers or other company employees, as opposed to the lab's analysts who would be operating the software in casework. For instance, the CBI v2.5 audit report begins with a memo from the lab's quality director stating: ``[ESR] submitted the summary and associated data files for the [audit] for the [laboratory]''~\cite[p.1]{cbiIV}. The company's copyright on the report further suggests developer involvement in the study. While the OCME v2.4 audit report does not explicitly state developer involvement, presence of the ESR copyright also suggests developer involvement~\cite{ocmeIV}. This evidence suggesting that developers may have conducted part or all of these audits raises concerns that even when tests incorporated user estimates, the users in the audit may not have been following the same protocols or training in the same environment as laboratory analysts who would be using the software in casework. 


\subsubsection{Gap: Inputs representative of casework samples vs. Inputs chosen out of convenience.}
\label{sec:scope-gap2-inputs}
\textbf{\\Goal: Compliant audits test the PGS system using inputs that represent the range of samples the lab (i) plans to analyze in casework and (ii) will likely encounter in casework.} ASB 018 requires that audits ``include case-type profiles''~\cite[R4.1.3]{asb018}, which the standard defines as samples ``exhibiting features that are representative of a plausible range of casework conditions [such as] [...] shared alleles''~\cite[p.1]{asb018}. ASB 018 also expects that these profiles ``represent [...] the range of actual casework samples intended for analysis with the system at the laboratory''~\cite[R4.1.3]{asb018}. 


\textbf{Practices:} \textbf{Audits primarily constructed test samples out of convenience---not out of a formal assessment of samples the lab intends to analyze or is likely to encounter in casework.} Labs frequently encounter mixtures containing DNA from related individuals in casework, and analysis of these mixtures is more likely to produce false positives due to individuals having similar DNA profiles (i.e., sharing alleles) \cite{kelly2022exploring}. While the CBI v2.5 and MSP v2.9.1 audits constructed mixtures specifically to test scenarios where contributors to the mixture are relatives~\cite{cbiIV, mspIV}, the PBSO v2.6.2 audit did not explicitly test system performance on mixtures with related contributors and instead described the mixtures in their test set as having ``varying amounts of allele sharing''~\cite[p.10]{pbsoIV}. The OCME v2.4 and v2.7 audit reports use similar language~\cite{ocmeIV, ocmeIV2.7}. This language suggest that these audits did not design their test set to account for mixtures of related individuals, especially since any set of mixtures is likely to have some amount of allele sharing (i.e., DNA profiles of any two unrelated individuals are likely to contain some of the same alleles). 


\subsubsection{How ASB 018's design enables these gaps}
\label{sec:scope-design}
\textbf{ASB 018's \textit{narrow definition} of the audit target misaligns with its vision that audits provide assurance of system reliability in casework.} ASB 018 acknowledges the importance of incorporating the software user's decisions into audits and implies a broad audit scope when describing audits as providing assurance of the reliability of the LR output in criminal cases. However, ASB 018 defines ``probabilistic genotyping system'' as the ``software, or software and hardware''~\cite[p.2]{asb018}, failing to acknowledge the lab analyst whose use of the PGS shapes the ultimate LR output. 

\textbf{The organization and language of ASB 018's requirements treat performance measurement and user inputs as separate items for audits to ``include.''} ASB 018 first requires that audits ``address [...] accuracy, sensitivity, specificity, and precision''~\cite[R4.1.3]{asb018}. In a separate requirement, ASB 018 requires that audits also ``include evaluati[ons of] user input parameters''~\cite[R4.1.4]{asb018}.  Instead of requiring that labs account for and measure how user inputs impact system accuracy, sensitivity, specificity, and precision, ASB 018 treats these two components as two separate criteria that audits can address in isolation of one another. For instance, labs can satisfy the first requirement by measuring PGS performance without considering user inputs and the second requirement with a separate study that changes user input parameters and observes whether LRs increase or decrease \textit{without} measuring performance~\cite{pbsoIV, mspIV}. 


\textbf{ASB 018 uses vague language that allows labs to do very little to comply.} As discussed earlier, ASB 018 requires that audits ``shall include evaluating user input parameters,'' but does not specify what this evaluation must entail~\cite[R4.1.4]{asb018}. ASB 018's requirement for labs to ``consider the effect of overestimating and underestimating the number of contributors'' is similarly vague. There is no requirement to measure the impact of user inputs on system accuracy, sensitivity, specificity, and precision. 


\textbf{ASB 018 envisions that labs conduct audits alone, but its use of vague language and failure to consider how developers might conduct part or all of an audit does not prohibit audits conducted by developers.} ASB 018 frames most requirements at the level of the internal validation study without specifying the actor responsible for carrying out the required action. When ASB 018 does address the lab, it does not clearly specify which audit components the lab is responsible for. Take, for example, the requirement that ``the laboratory shall validate [the] PGS system''~\cite[R4.1]{asb018}. Does a lab's review of a summary of the audit conducted by the developers count as the lab ``validating'' the software? Can the developers provide the inputs to the software using their estimates for user inputs such as the number of contributors? 


\textbf{ASB 018 does not require that labs take note of the scenarios they typically encounter or intend to analyze---a preliminary step that would hold labs accountable to some specific range of samples.} ASB 018 requires that audits test ``case-type profiles [...] that represent [...] the range of actual casework samples intended for analysis''~\cite[R4.1.3]{asb018}, but does not require labs to determine and specify this range---a preliminary step that would help hold labs accountable for creating test sets that represent this range of intended uses. 



\subsection{Standards Identification: Articulating the Criteria or Expectations the System is Held to in the Audit}
\label{sec:stds}

\subsubsection{Gap: Predetermine specifications that define acceptable system performance vs. Failure to explicitly define specifications}
\label{sec:stds-gap1}
\textbf{\\Goal: Labs test the PGS system against predetermined specifications: expectations for acceptable system performance.} In its definition of ``validation,'' ASB 018 emphasizes the importance of ensuring the PGS system ``will consistently [meet] its predetermined specifications''~\cite[p.2]{asb018}. Taken together with Requirement 4.1.3 and ASB 018's vision for audit purpose, we understand ASB 018 to further suggest that these specifications must define acceptable ``accuracy, sensitivity, specificity, and precision'' of the PGS system's LR output~\cite[R4.1.3]{asb018}. 


\textbf{Practices: Audit reports often do not explicitly define expectations for acceptable software behavior, suggesting that audits were not evaluating the system against predetermined expectations for acceptable performance.} Instead of explicitly stating the expectations they were evaluating against, audits sometimes implicitly communicated their expectations by describing observed results as ``expected.'' For example, after presenting observed true positive rates, the MSP v2.9.1 audit report concludes: ``As expected, there was a high percentage of [true positive] LRs''~\cite[p.27]{mspIV}. This language suggests that the lab may have understood a ``high'' true positive rate to be acceptable system sensitivity, but the audit never explicitly stated this expectation. At other times, audits simply concluded by presenting results without any language suggesting that results met expectations, making it difficult to infer what the lab's expectations were or whether they had \textit{any} expectations to begin with. For example, the OCME v2.4 audit measured software precision (i.e., the degree of variability in LRs between multiple software runs on the same inputs) and concluded that the observed LRs ``demonstrate that [...] the values obtained for the various runs remain close; within one order of magnitude except for the four-person mixture''~\cite[p.19]{ocmeIV}. 



\subsubsection{Gap: Expectations align with forensic labs' customers' expectations for PGS performance vs. Expectations that may conflict with customers' desires}
\label{sec:stds-gap2}
\textbf{\\Goal: Compliant audits test the system against expectations for acceptable system performance that align with lab customers' perceptions of what constitutes acceptable system performance.} ASB 018 emphasizes that ``[audits] provide the study results and conclusions necessary for customers of [forensic labs] to have confidence in the [LR] evidence provided''~\cite[Foreword]{asb018}. In order for audits to support customer\footnote{We interpret ASB 018's use of the word ``customer'' to imply stakeholders who pay laboratories for their services (e.g., prosecutors). Importantly, ASB 018 does not use the word ``stakeholder'', which refers to a broader set of actors.} assurance of PGS system reliability, the audit must assess PGS system behavior against lab customers'  expectations (as opposed to the PGS developer or forensic lab's expectations).


\textbf{Practices:} \textbf{Audit reports suggest that audits relied on expectations for acceptable system performance that may misalign with customers' needs, and none of the audit reports mention customers' needs or perceptions.} For example, the MSP v2.9.1 audit report concludes that the audit's observed false positive rates of up to 1.3\% and false negative rates of up to 16\% demonstrate that the PGS performed ``as expected''~\cite[p.27]{mspIV}. In doing so, the audit did not consider how a defense attorney's or a prosecutor's perception of acceptable false positive or false negative rates may differ. Furthermore, we found that none of the audits even mentioned customers' needs or perceptions in their reports, suggesting that labs did not consider whether their expectations aligned with those of their customers.


\subsubsection{Gap: Falsifiable expectations vs. Subjective, otherwise unclear, and contradictory expectations.}
\label{sec:stds-gap3}
\textbf{\\Goal: } \textbf{Expectations for acceptable system performance are falsifiable.} In other words, labs define what constitutes acceptable system behavior using criteria that enable labs to assess observed system behavior against these criteria and unambiguously state whether the system meets these criteria. Crucially, labs must be able to unambiguously identify when software behavior \textit{fails} to meet these criteria for acceptable performance, since the end goal of an audit is to establish boundaries on acceptable software use.


\textbf{Practices: }\textbf{Audits described expectations using ambiguous criteria.} For instance, audits frequently defined expectations using subjective criteria such as ``high'' and ``low.'' For example, all audits stated that the LR for true contributors ``should be high.'' The MSP v2.9.1 audit report concludes that their observed sensitivity rates demonstrated ``high [PGS] sensitivity''~\cite[p.23]{mspIV}. Audit reports also use other subjective criteria when discussing expectations, such as ``intuitive'' (e.g., ``result in an intuitive inclusionary LR''~\cite[p.7]{ocmeIV2.7}), ``generally'' (e.g., ``true contributors generally gave high LRs''~\cite[p.21]{ocmeIV2.7}), and ``close'' (e.g., ``the values obtained for the various runs remain close''~\cite[p.19]{ocmeIV}). \textbf{Audits also used otherwise unclear language in their expectations.} All audits' definitions of sensitivity imply that they expected the PGS to ``reliably resolve the DNA profile of [true] contributors''~\cite{ocmeIV, ocmeIV2.7, pbsoIV, mspIV, cbiIV}. However, none of the audit reports specify what it means for a PGS to ``resolve'' a DNA profile. For example, does an LR above 1 count as ``resolving'' the profile, or should the LR be above some other threshold value like 1,000? \textbf{Expectations within the same report were sometimes contradictory.} For example, the OCME v2.4 audit report defines specificity as ``the ability of the software to reliably exclude non-contributors''~\cite[p.6]{ocmeIV}. However, when discussing experiment results, the report states that some false positive LRs were ``expected''~\cite[p.14]{ocmeIV}. 


\subsubsection{How ASB 018's design enables these gaps}
\label{sec:stds-design}
While ASB 018 refers to ``predetermined specifications'' in its definition of validation, \textbf{the standard does not include any requirement that labs actually \textit{define} these specifications}. Furthermore, despite emphasizing the importance of audits establishing boundaries on acceptable software use, \textbf{ASB 018 does not establish any requirements for \textit{how} these specifications should be defined to ensure that defined expectations actually support the end goal to establish boundaries on acceptable software use}. Instead, ASB 018 simply requires that labs ``\textit{address} [...] accuracy, sensitivity, specificity, and precision''~\cite[emph. added]{asb018}. Similar to our discussion of ASB 018's use of the word ``consider'' in Section~\ref{sec:scope-design}, \textbf{the word ``address'' allows labs to comply without clearly and explicitly defining what constitutes acceptable accuracy, sensitivity, specificity, and precision and then assessing PGS system performance against these predetermined specifications}. 



\subsection{Performance Analysis: The Actual Evaluation Itself (Gathering Data on System Performance)}
\label{sec:performance}

\subsubsection{Gap: Measure system performance on a variety of sample types vs. Measure performance on a narrow subset of samples tested.} 
\label{sec:performance-gap1}
\textbf{\\Goal: Labs measure sensitivity, specificity, precision of the PGS system's LR outputs for a variety of sample types, and labs measure accuracy for a smaller set of sample types.}
A goal that follows from \textit{defining} expectations for acceptable PGS system performance with respect to sensitivity, specificity, precision, and accuracy is \textit{measuring} these dimensions of system performance. Furthermore, ASB 018 envisions that labs generate these measurements for a variety of DNA sample types---a detail crucial to ensuring that audits help establish boundaries on software use. ASB 018 also envisions that labs measure system accuracy for a variety of sample types but explicitly limits its expectations for labs to a smaller set of sample types: ``[one-person] samples [and simple][two]-person mixtures''~\cite[p.1]{asb018}. ASB 018 emphasizes that situations where ``the ground truth is not known are not suitable for accuracy studies''~\cite[p.1]{asb018}


\textbf{Practices: Audits often measured LR sensitivity, specificity, precision, and accuracy on a narrow subset of the samples that labs had access to or could have tested.} Every report documents audit measurements of software sensitivity and specificity in an experiment specifically dedicated to sensitivity and specificity. \textbf{In these experiments, most audits did \textit{not} explicitly test DNA sample types that labs often encounter in casework, such as significantly degraded\footnote{E.g., DNA samples degrade after exposure to sunlight for several days} mixtures}. Instead, reports only address degraded samples when describing later experiments that often did not measure system specificity, precision, and accuracy. For example, the OCME v2.4 audit artificially degraded a single one-person DNA sample and only investigated the impact of degradation on the true contributor's LR (instead of also investigating non-contributor LRs to test specificity)~\cite{ocmeIV}. \textbf{All audits investigated system precision in separate experiments, and most audits' precision experiment test sets were  \textit{significantly smaller subsets} of their sensitivity and specificity experiment test sets.} Lastly, while ASB 018 explicitly constrains the set of sample types for which labs should measure system accuracy, ASB 018 \textit{does} include simple two-person mixtures in its set of suitable sample types. However, \textbf{all audits' accuracy experiments only evaluated one-person samples.} Audits sometimes even acknowledged that the samples they were testing were a subset of the entire set of sample types they could have tested, stating something like: ``There is a small subset of profiles where [system accuracy can be tested]. These include [one-person samples]''~\cite[p.3]{ocmeIV}. In other words, these audits measured accuracy on one type of sample (which is also one of the types of DNA samples least likely to be associated with poor PGS performance) when they could have also included two-person mixtures.


\subsubsection{Gap: Robust audit measurements vs. Inadequate consideration of sample size sufficiency} 
\label{sec:performance-gap2}
\textbf{\\Goal: Labs produce \textit{robust} measurements of system performance through repeated testing to account for variability in measurements.} ASB 018 recognizes that performance measurements may vary due to upstream DNA preparation procedures, stochasticity in the software's underlying statistical models, and differences between DNA samples. As such, the standard highlights the importance of ``repeated testing'' on a ``sufficient'' range of sample types and number of samples of each type~\cite[Annex A, R4.1.3]{asb018}. 


\textbf{Practices:} \textbf{Audits neither commented on the sufficiency of the sample sizes used to test system performance, nor provided measures used to determine whether sample sizes were sufficient.} This lack of discussion is especially concerning since sample sizes varied drastically between audits. For example, when testing system precision, the MSP v2.9.1 audit used one four-person sample~\cite[p.65]{mspIV}, whereas the OCME v2.7 audit used eight samples spanning two- to five- contributor mixtures~\cite[p.28]{ocmeIV2.7}. For sensitivity and specificity, while the MSP v2.9.1 audit tested 30 two-person and 17 five-person mixtures~\cite[p.14-15]{mspIV}, the OCME v2.7 audit tested 169 two-person and 40 five-person mixtures~\cite[p.15]{ocmeIV2.7}. 


\subsubsection{How ASB 018's design enables these gaps}
\label{sec:performance-design}
\textbf{ASB 018's vague language and use of separate line-level requirements to address interrelated concepts allows labs to analyze PGS performance on narrow subsets of test sets.} While ASB 018 requires that internal validation studies to address ``accuracy, sensitivity, specificity, and precision'' in one line-level requirement, another line-level requirement states, ``These studies shall include [samples] that represent [....] the range of actual casework samples intended for analysis''~\cite[R4.1.3]{asb018}. In this second requirement, ASB 018 does not specify whether ``[t]hese studies'' broadly refers to audits or narrowly refers to individual accuracy, sensitivity, specificity, and precision experiments within these audits (which the standard also refers to as ``studies'' elsewhere, e.g., ``accuracy studies''~\cite[p.1]{asb018}). A permissive interpretation of the second requirement would interpret ``[t]hese studies'' to mean ``[i]nternal validation studies,'' which would allow for a precision experiment to test a narrow subset of test samples while the audit as a whole satisfies the requirement for audits to test representative DNA samples. Furthermore, as we have discussed earlier (Section~\ref{sec:scope-design}), ASB 018 treats complex sample characteristics such as degradation and software performance measurement as two separate criteria that audits can address in isolation from one another. 


\textbf{ASB 018 lists only three dimensions of DNA sample types in the requirements, failing to capture key dimensions that the standard outlines in its Terms and Definitions section.} The standard requires that audits include case-type profiles ``that represent \textit{(in terms of number of contributors, mixture ratios, and total DNA [amount])} the range of actual casework samples intended for analysis''~\cite[R4.1.3, emph. added]{asb018}. Here, the standard names three DNA sample characteristics that labs must account for: the number of contributors, mixture ratios, and the total DNA amount. However, the standard's definition of ``case-type profiles'' names other dimensions to consider, such as the amount of ``degradation'' and the number of ``shared alleles'' between individual contributors in a mixture~\cite[p.1]{asb018}. ASB 018's omission of key factors like degradation and allele sharing not only limits attention to the dimensions specified in the requirements, but also introduces ambiguity around whether the dimensions mentioned in the Terms and Definitions section are requirements or suggestions. 
 

\textbf{ASB 018 neither provides criteria for unambiguously determining whether the sample size of a lab's test set is sufficient, nor requires labs to establish their own criteria.} Instead of specifying a minimum number that labs must comply with, ASB 018 requires that labs ``shall perform sufficient studies''~\cite[Annex A, R4.1.3]{asb018}. However, the standard neither specifies how labs should operationalize ``sufficient,'' nor requires labs to establish their own criteria for determining sufficiency. While the standard's requirements language provides labs with the flexibility to design studies that fit their needs and contexts, it also enables labs to test any types and number of samples (e.g., two samples) without justifying why their choice is sufficient. In other words, the standard does not hold labs accountable to any definition of ``sufficient.''


\subsection{Post-Audit Judgment: Discussing Results and Translating Results Into Actions or Consequences}
\label{sec:judgement}

\subsubsection{Gap: Conclude whether measured behavior is acceptable vs. Frame software failures as correct behavior.} 
\label{sec:judgement-gap1}

\textbf{\\Goal: Labs assess observed behavior and determine whether observed behavior is acceptable or unacceptable.} \textit{Providing judgment} about observed system behaviors, rather than stopping at simply producing these observations (i.e., measuring system behaviors), is crucial to ensuring that audits establish boundaries on acceptable software use in casework. In other words, the lab must label a test result as acceptable or unacceptable, instead of simply stating the test result. 


\textbf{Practices:  Audits frequently framed software failures as correct behavior. Most audits did not use the word ``failure'' to describe incorrect results, false positives, and false negatives.} For instance, the OCME v2.4 audit report frequently reframes observed false positive results as correct behavior. The audit report states that the false positives were ``not a failure of STRmix''~\cite[p. 13]{ocmeIV}. The OCME v2.7 audit report uses similar language~\cite[p. 24]{ocmeIV2.7} and also frames other software failures as ``non-intuitive,'' stating that some test runs ``resulted in non-intuitive results with an LR = 0 or failure to run''~\cite[p.14]{ocmeIV2.7}. Audits also frequently framed software failures as ``not unexpected'' instead of as failures, responding to false positive and false negative LRs with statements like ``these results are not unexpected''~\cite[p.24]{ocmeIV2.7}. None of the audit reports analyzed in our study frame observed failures as ``unacceptable.'' 


\subsubsection{Gap: Establish boundaries on acceptable software use in casework vs. Not establishing such boundaries}
\label{sec:judgement-gap2}
\textbf{\\Goal: Labs use their judgments of observed system behavior to establish boundaries on acceptable software use in casework.} As we described at the beginning of Section~\ref{sec:results}, ASB 018 envisions that audits ultimately establish boundaries on acceptable software use, and these boundaries must be specified in terms of DNA sample types.  
  

\textbf{Practices: }\textbf{None of the audits used observed software failures to define boundaries on software use in terms of DNA sample types.} Some audits established another type of boundary: an ``uninformative'' range for LRs (e.g., 0.001 - 1000~\cite[p.24]{ocmeIV2.7}) that the lab would use in criminal cases to present LRs falling in this range as ``inconclusive'', i.e., neither inculpatory nor exculpatory. However, this boundary is not a \textit{restriction} on software use and does not prevent labs from introducing false positive LRs as evidence in criminal cases when the PGS system produces extreme false positive LRs (which is common in the case of marginal samples where there is a high degree of allele sharing between a true contributor and a non-contributor).  


\textbf{Oftentimes, in place of boundaries on software use, audits emphasized that analysts would be able to identify the issue and adjust their software use accordingly}. For example, upon observing false positive results, one audit concluded that such scenarios are ``recognizable by trained analysts'' who would be able to correct the analysis accordingly to avoid such errors~\cite[p.8]{ocmeIV2.7}. However, as we discussed in Section~\ref{sec:scope-gap1-target}, audits were scoped around the PGS instead of the broader sociotechnical system, so these audits do not actually provide assurance that analysts would be able to avoid such errors in casework. \textbf{Audits also sometimes concluded without providing any assurance of reliability in casework, instead advising analysts to proceed with caution.} For example, after observing several false negative results, the CBI v2.5 audit simply advised analysts to ``exhibit caution [...] for situations such as this''~\cite[p.31]{cbiIV}. 

\subsubsection{How ASB 018's design enables these gaps}
\label{sec:judgement-design}
\textbf{ASB 018 neither clearly defines ``specification'', nor requires labs to explicitly label whether observed performance is acceptable.} In Section~\ref{sec:judgement-gap1}, we discussed how labs framed software failures as ``not unexpected,'' instead of labeling these failures as unacceptable behavior that the lab should prevent in casework through establishing boundaries on acceptable software use based on audit results. These practices further reveal consequences of the standard's failure to specify what constitutes a ``specification'' (see Section~\ref{sec:stds-design}), and also reveal that the standard fails to ensure that labs actually judge whether audit results are acceptable---a crucial step towards establishing boundaries on acceptable software use.


\textbf{ASB 018 does not explicitly require that audits establish boundaries on software use. Instead, ASB 018 relies on the word ``limitations'', which the standard does not define.} Despite repeatedly describing that audits help determine limitations, ASB 018 does not define what constitutes a limitation. The accompanying factsheet clarifies that limitations are boundaries that define ``the range of DNA profiles upon which the program may be used effectively''~\cite{asb018factsheet}. However, these details are crucially lacking from the standard, the sole document that studies would be assessed against for compliance. Furthermore, \textbf{despite emphasizing that audits establish ``limitations'', ASB 018 does not require a lab to actually \textit{specify} the limitations they identify in the audit.} In other words, the standard seeks to guarantee that labs use audits to determine system limitations, but makes no requirement that labs actually do so. 




\subsection{Audit Communication: Documenting Audit Design, Methodology, and Results} 
\label{sec:communication}


\subsubsection{Gap: Documentation enables third party scrutiny and reproducibility vs. Documentation contains insufficient information about test samples and results}
\label{sec:communication-design-gap}\textbf{ \\Goal: Lab documentation of audit data, methods, and results enables third party scrutiny.} 
ASB 018 emphasizes the importance of audit documentation that addresses the information needs of third parties, stating: ``it is incumbent upon any laboratory performing these [audits] to retain these results for the examination and evaluation by third parties. The results should be documented in such a way that the [...] [audits] can be reproduced and decisions made on the basis of these studies documented''~\cite[p.5]{asb018}.

\textbf{Practices: Audit reports often do not provide sufficient information about test samples and conditions that third parties could use to interpret study results.} The audit reports we analyzed omit crucial information about the range of DNA sample types examined, often describing test samples using vague language. For example, each audit report describes their sensitivity and specificity experiment test samples as ``varying'' in their ``[total] DNA [amounts,] mixture proportions'', and ``amounts of allele sharing'' (e.g.,~\cite[p.10]{pbsoIV}). When audit reports do provide specific ranges of sample types explored, they often only provide details for a small set of DNA sample characteristics (e.g., specifying the range of total DNA amounts studied, but not specifying the range of mixture proportions studied~\cite[p.15]{ocmeIV2.7}). \textbf{Audit reports also often do not provide data that would enable third party verification of their claims.} For example, the OCME v2.4 audit report claims that ``The results of all comparisons was expected. [...] the [non-contributors got] very low or 0 LRs''~\cite[p.22]{ocmeIV}. In the appendix, the audit report documents the number of non-contributors with LRs = 0 and plots the LRs in scatterplots, but does not provide exact LR values for non-contributors, despite the plots revealing several false positive LRs~\cite[p.47-50]{ocmeIV}.



\subsubsection{How ASB 018's design enables these gaps}
\label{sec:communication-design}
\textbf{ASB 018 requires that labs document their audits, but non-mandatory wording to describe the importance of ensuring that documentation is accessible and useful to third parties}. ASB 018 requires documentation by stating that ``All [audits] shall be documented and retained by the laboratory'', but switches to using the phrase ``it is incumbent upon'' when emphasizing that labs should document the results to support ``examination and evaluation by third parties''~\cite[Annex A R4.5]{asb018}. \textbf{ASB 018 also fails to specify specific pieces of information that labs must document,} even though the standard specifies in the line immediately after that ``Laboratories shall have a \textbf{summary statement of the sample types} of which the developer used to for their developmental validation''~\cite[Annex A R4.5]{asb018} While the standard could specify a similar requirement that labs document the sample types they tested in their audit, the standard does not include any documentation requirement with this level of specificity. 
